\section{Lei de Faraday}

A lei da indução de Faraday relaciona o campo elétrico com a variação no campo magnético. Experiências mostram que ao aproximar ou afastar um imâ de uma espira condutora, uma corrente elétrica surge nela. 

Seja $C$ um circuito fechado. Se $\Phi_B$ é o fluxo do campo magnético através de uma superfície $S$ delimitada por $C$, então temos
\begin{equation}
\mathcal{E} = -\dot{\Phi}_B,
\label{eq: lenz}
\end{equation}

\noindent onde $\mathcal{E}$ é a força eletromotriz equivalente à corrente gerada no condutor. A eq.\eqref{eq: lenz} é a lei de Lenz, e diz que variações no fluxo magnético através de um condutor geram correntes que se opõem a essa variação. Por exemplo, ao aproximar um imã de uma espira condutora, a corrente gerada irá, por sua vez, gerar um novo campo magnético, que causará uma força de repulsão no imã. Ao afastar o imã, o mesmo ocorre, mas a força é de atração. Em ambos os casos, a força gerada pela corrente se opões a variação de fluxo magnético. 

Decorre da eq.\eqref{eq: int_gauss_mag} que o fluxo de $\mathbf{B}$ sobre qualquer superfície não fechada depende apenas de seu contorno. Portanto a eq.\eqref{eq: lenz} é valida para qualquer superfície de contorno $C$. Como a f.e.m. é a circulação do campo elétrico sobre um circuito, temos que 
\begin{equation}
\oint_C \mathbf{E} \cdot d\mathbf{l} = -\frac{d}{dt} \iint_S \mathbf{B} \cdot d\mathbf{S}.
\end{equation}

\noindent Utilizando, o teorema do rotacional (\ref{th: stokes}), obtemos 
\begin{equation}
\nabla \times \mathbf{E} = - \partial_t \mathbf{B}.
\label{eq: faraday}
\end{equation}

Embora a teoria tenha sido formulada no contexto de correntes elétricas num circuito, as equações mostradas nessa sessão são válidas mesmo na ausência de condutores. 

